% Exercise 2.13
\begin{itemize}
    \item[(a)] Suppose we have $x_0 < y_0$ and
        $\sum_{i=0}^n{x_i} > \sum_{i=0}^n{y_i}$, then we have
        \[
            \sum_{i=0}^n{x_i - y_i} > 0,
        \]
        and, since $x_0 < y_0 \Rightarrow x_0 - y_0 < 0$, we have
        \[
            \sum_{i=1}^n{x_i - y_i} > x_0 - y_0.
        \]
        Letting $a_i = x_i - y_i$ for $i = 0, 1, \ldots, n$, we have $a_i < 0$
        and $\sum_{i=1}^n{a_i} > 0$.

        Now, we have $\text{PV}_x(d) = \sum_{i=0}^n{x_i d^i}$ and
        $\text{PV}_y(d) = \sum_{i=0}^n{y_i d^i}$, so that
        $\text{PV}_a(d) = \sum_{i=0}^n{a_i d^i}$. We then have $a_i < 0$ and
        $\sum_{i=1}^n a_i > 0$, so that $\text{PV}_a(0) = 0$ and
        $\text{PV}_a(d_1) > 0$ for some $d_1 > 0$.

        Since $\text{PV}_a(d)$ is a polynomial function with independent
        variable $d$, and is therefore continuous throughout its domain, by the
        Intermediate Value Theorem, we have $\text{PV}_a(c) = 0$ for some
        $c > 0$ s.t. $c < d_1$, so that $\text{PV}_x(c) = \text{PV}_y(c)$.
    \item[(b)] We have $x_0 = -100$ and $y_0 = -150$, and
        $\sum_{k=0}^5{x_k} = 50$, while $\sum_{k=0}^5{y_k} = 60$, so that the
        cash flow streams fit our assumptions, for the roles of $x_i$ and $y_i$
        swapped. Thus we must have some $c > 0$ s.t.
        $\text{PV}_x(c) = \text{PV}_y(c)$, or, $\text{PV}_a(c) = 0$, as in (a),
        but with $a_k = y_k - x_k$.

        Thus we have $a_0 = -50$ and
        \begin{equation*}
            \begin{split}
                \text{PV}_a(c) &= \sum_{k=0}^5 a_k c^k \\
                &= -50 + \sum_{k=1}^5 (42 - 30) c^k \\
                &= -50 + \sum_{k=1}^5 12 c^k.
            \end{split}
        \end{equation*}
        Letting $c = 1 / (1 + r)$ and using Newton's method, as in Exercise 11,
        to solve for the single positive rate $r$ at which
        $\text{PV}_a(c) = 0$, we have the following:
        \begin{lstlisting}
>>> from investment_science.cash_flows import irr_iterates, present_value
>>> from investment_science.ch02.ex2_13 import DIFFERENCE
>>> DIFFERENCE
[-50.0, 12.0, 12.0, 12.0, 12.0, 12.0]
>>> rs = irr_iterates(DIFFERENCE, 0.1, 3)
>>> for r in rs:
...     print(f"r = {r}, f(r) = {present_value(DIFFERENCE, r)}")
...
r = 0.1, f(r) = -4.5105587670986305
r = 0.061186157304593224, f(r) = 0.3856708315489712
r = 0.0640051684113512, f(r) = 0.0023300017467331458
r = 0.06402240700516368, f(r) = 8.621847769063606e-08
>>> 1 / (1 + rs[-1])
0.9398298319812988
        \end{lstlisting}
        Thus we have $r \approx 0.06402240700516368$, so that the crossover
        value is given by $c \approx 0.9398298319812988$, i.e., the present
        values of the two projects are equal at an interest rate of
        approximately 6.40\%.
\end{itemize}
